\itemhead{复习题一 · 第15题（教材p.30）}{XJ-B1-C01-RV-P030-Y0883}
\label{XJ-B1-C01-RV-P030-Y0883}
{\footnotesize 湘教版教材p.30 / PDF第37页 / 复习题一 / 第15题}
\field{原题}{为完成一项实地测量任务，夏令营的同学们成立了一支“测绘队”，需要 24 人参加测量，20 人参加计算，16 人参加绘图。测绘队的成员中很多同学是多面手，已知有 8 人只参加测量、计算，有 6 人只参加测量、绘图，有 4 人只参加计算、绘图，另有几人三项工作都参加了。试问这支测绘队至少有多少人？}
\field{参考答案}{至少30人。设三项都参加者为 $t$，则三类只参加一项的人数为 $10-t,8-t,6-t$，故 $0\le t\le6$；总人数为 $42-2t$，取 $t=6$ 得30。}
{\small \field{关联知识}{S06 Venn图、容斥与实际计数（教材帮p.22 / PDF31）}
\field{方法与易错点}{先画互不重叠区域，从最内层填写人数，再依据总人数或非负限制求解。 易错：把只参加两项的人数当作完整交集；重复计算三重交集；忽略一种都不参加。}
\field{答案依据}{答案册相关题组：PDF第8、9页（册内页7、8）；按当前题库原题逐项重算，不按旧版题号机械对套。}}
